% Ad Quintum Fratrem 2.1 (ad-quintum-fratrem-02-01)
% Source: https://raw.githubusercontent.com/PerseusDL/canonical-latinLit/master/data/phi0474/phi058/phi0474.phi058.perseus-lat1.xml
% Fetched by scripts/fetch_latin.py

% Dateline (Perseus): Scr. Romae m. Dec. a. 697 (57) .

\ciceroLetterOpener{MARCVS QVINTO FRATRI SALVTEM}

\ciceroSection{1} epistulam quam legisti mane dederam sed fecit i humaniter Licinius quod ad me misso senatu vesperi venit, ut si quid esset actum ad te, si mihi videretur, perscriberem. senatus fuit frequentior quam putabamus esse posse mense Decembri sub dies festos. consulares nos fuimus et duo consules designati, P. Servilius, M. Lucullus, Lepidus, Volcacius, Glabrio, praetores sane frequentes fuimus, omnino ad Cc. commorat exspectationem Lupus ; egit causam agri Campani sane accurate. auditus est magno silentio. materiam rei non ignoras. nihil ex nostris actionibus praetermisit. fuerunt non nulli aculei in Caesarem, contumeliae in Gellium, expostulationes cum absente Pompeio. causa sero perorata sententias se rogaturum negavit, ne quod onus simultatis nobis imponeret; ex superiorum temporum conviciis et ex praesenti silentio quid senatus sentiret se intellegere dixit Milo. coepit dimittere. tum Marcellinus ' noli ,' inquit, 'ex taciturnitate nostra, Lupe, quid aut probemus hoc tempore aut improbemus iudicare. ego , quod ad me attinet, itemque arbitror ceteros, idcirco taceo quod non existimo, cum Pompeius absit, causam agri Campani agi convenire.' tum ille se senatum negavit tenere.

\ciceroSection{2} Racilius surrexit et de iudiciis referre coepit. Marcellinum quidem primum rogavit. is cum graviter de Clodianis incendiis, trucidationibus, lapidationibus questus esset, sententiam dixit, ut ipse iudices per praetorem urbanum sortiretur, iudicum sortitione facta comitia haberentur ; qui iudicia impedisset, eum contra rem publicam esse facturum. adprobata valde sententia C. Cato contra dixit et C. Cassius maxima acclamatione senatus, cum comitia iudiciis anteferrent. Philippus adsensit Lentulo. postea Racilius de privatis me primum sententiam rogavit.

\ciceroSection{3} multa feci verba de toto furore latrocinioque P. Clodi ; tamquam reum accusavi multis et secundis admurmurationibus cuncti senatus. orationem meam conlaudavit satis multis verbis non me hercule indiserte vetus Antistius, isque iudiciorum causam suscepit antiquissimamque se habiturum dixit. Ibatur in eam sententiam. tum Clodius rogatus diem dicendo eximere coepit ; furebat a Racilio se contumaciter urbaneque vexatum. deinde eius operae repente a Graecostasi et gradibus clamorem satis magnum sustulerunt, opinor, in Q. Sextilium et amicos Milonis incitatae. eo metu iniecto repente magna querimonia omnium discessimus. habes acta unius diei ; reliqua, ut arbitror, in mensem Ianuarium reicientur. de tribunis pl. longe optimum Racilium habemus. videtur etiam Antistius amicus nobis fore ; nam Plancius totus noster est. fac , si me amas, ut considerate diligenterque naviges de mense Decembri.
